% ─────────────────────────────────────────────────────────────────────────────
% CAPÍTULO 1 — Arquitectura General y Topología de Red
% ─────────────────────────────────────────────────────────────────────────────
\chapter{Arquitectura General del Sistema}

\section{Resumen Ejecutivo}

El \textbf{Hermes Stack} es una plataforma de agente conversacional IA
completamente autohospedada, desplegada en un único VPS y accesible
públicamente a través de subdominios HTTPS con certificados renovados
automáticamente.

\begin{infobox}[Propuesta de valor clave]
Soberanía total sobre los datos · Coste mensual fijo predecible ·
Latencia end-to-end de voz menor a 500 ms · Zero-downtime deployment
automatizado desde GitHub.
\end{infobox}

\subsection{Métricas objetivo}

\begin{table}[h]
\centering
\caption{KPIs del Hermes Stack}
\begin{tabularx}{\textwidth}{lXr}
\toprule
\textbf{Métrica} & \textbf{Descripción} & \textbf{Objetivo} \\
\midrule
Latencia E2E voz & Audio entrada → Audio salida & $<$ 500 ms \\
Uptime mensual   & Disponibilidad del servicio Hermes & $\geq$ 99.5\% \\
TTFB LLM         & Primer token desde LiteLLM & $<$ 150 ms \\
Health check     & Ciclo de auto-curación Autoheal & 30 s \\
Deploy time      & Push a main → servicio actualizado & $<$ 3 min \\
\bottomrule
\end{tabularx}
\end{table}

\section{Diagrama General del Sistema}

El diagrama siguiente muestra los tres planos del sistema: el plano de
\textit{usuario} (clientes externos), el plano de \textit{control}
(Caddy + GitHub Actions) y el plano de \textit{datos} (servicios internos).

\begin{figure}[h]
\centering
\begin{tikzpicture}[
  node distance=0.8cm and 1.4cm,
  svc/.style={draw=hermesblue,fill=hermeslight,rounded corners=6pt,
              minimum width=2.4cm,minimum height=0.9cm,
              font=\sffamily\small\bfseries,align=center},
  ext/.style={draw=hermesaccent,fill=hermesaccent!10,rounded corners=6pt,
              minimum width=2.4cm,minimum height=0.9cm,
              font=\sffamily\small,align=center},
  cloud/.style={draw=gray,fill=gray!10,ellipse,
                minimum width=2.6cm,minimum height=1cm,
                font=\sffamily\small,align=center},
  arr/.style={->,>=Stealth,thick,hermesblue},
  label/.style={font=\sffamily\tiny,gray}
]

%% — Capa usuario
\node[cloud] (user)  {Navegador /\\App Android};
\node[cloud,right=2cm of user] (gh) {GitHub\\Actions};

%% — Caddy
\node[svc,below=1.2cm of user,fill=hermesblue!30] (caddy)
  {Caddy\\\small HTTPS + TLS};

%% — Servicios internos (fila 1)
\node[svc,below left=1.2cm and 0.6cm of caddy]  (hermes)   {Hermes\\Agent :8080};
\node[svc,below=1.2cm of caddy]                 (litellm)  {LiteLLM\\ :4000};
\node[svc,below right=1.2cm and 0.6cm of caddy] (grafana)  {Grafana\\ :3000};
\node[svc,right=0.5cm of grafana]               (website)  {Website\\ :3001};

%% — Servicios internos (fila 2)
\node[svc,below=1.0cm of hermes]  (whisper) {Whisper\\STT :9000};
\node[svc,below=1.0cm of litellm] (redis)   {Redis\\Cache};
\node[svc,below=1.0cm of grafana] (prom)    {Prometheus\\ :9090};

%% — APIs externas
\node[ext,right=1.5cm of litellm] (openrouter) {OpenRouter\\/ Gemini};
\node[ext,below right=0.4cm and 1.5cm of hermes] (elevenlabs) {ElevenLabs\\Flash v2.5};

%% — Flechas
\draw[arr] (user)    -- node[label,right]{HTTPS} (caddy);
\draw[arr] (caddy)   -- (hermes);
\draw[arr] (caddy)   -- (litellm);
\draw[arr] (caddy)   -- (grafana);
\draw[arr] (caddy)   -- (website);
\draw[arr] (hermes)  -- node[label,left]{STT} (whisper);
\draw[arr] (hermes)  -- (litellm);
\draw[arr] (hermes)  -- node[label,right]{TTS WS} (elevenlabs);
\draw[arr] (litellm) -- (redis);
\draw[arr] (litellm) -- node[label,above]{LLM API} (openrouter);
\draw[arr] (prom)    edge[bend left=20] (hermes);
\draw[arr,dashed,gray] (gh) -- node[label,right]{SSH deploy} (caddy);

\end{tikzpicture}
\caption{Arquitectura general del Hermes Stack}
\end{figure}

\section{Principios Arquitectónicos}

\begin{enumerate}[label=\textbf{\arabic*.}]
  \item \textbf{Soberanía de datos:} Todos los datos de usuario transitan
    por el VPS propio. Sólo salen peticiones a APIs externas (LLM, TTS).

  \item \textbf{Red en capas:} Dos redes Docker aisladas (\code{backend},
    \code{monitoring}). Ningún contenedor de datos está expuesto al host.

  \item \textbf{Secretos en fichero .env:} El servidor jamás expone claves
    en variables de entorno del sistema operativo ni en el repositorio.

  \item \textbf{Inmutabilidad de contenedores:} Todos los servicios tienen
    \code{restart: unless-stopped} y healthchecks declarativos.

  \item \textbf{Auto-curación:} Autoheal monitoriza cada 30 s y reinicia
    cualquier contenedor con healthcheck fallido.
\end{enumerate}

\section{Topología de Red Docker}

El stack utiliza dos redes bridge internas:

\begin{table}[h]
\centering
\caption{Redes Docker del Hermes Stack}
\begin{tabularx}{\textwidth}{llX}
\toprule
\textbf{Red} & \textbf{Driver} & \textbf{Miembros} \\
\midrule
\code{backend}    & bridge & litellm, redis, whisper-stt, hermes, website, autoheal, filebrowser \\
\code{monitoring} & bridge & prometheus, grafana, cadvisor, node-exporter, hermes \\
\bottomrule
\end{tabularx}
\end{table}

\begin{infobox}[Aislamiento de red]
Los contenedores de \textit{monitoring} no pueden iniciar conexiones
hacia los contenedores de \textit{backend} (salvo Hermes y LiteLLM que
pertenecen a ambas redes para exponer métricas a Prometheus).
\end{infobox}

\subsection{Puertos expuestos al host}

Todos los puertos están ligados a \code{127.0.0.1} (loopback). Caddy
actúa como único punto de entrada público en los puertos 80/443.

\begin{table}[h]
\centering
\caption{Mapa de puertos del stack}
\begin{tabularx}{\textwidth}{llXl}
\toprule
\textbf{Servicio}  & \textbf{Puerto host} & \textbf{Subdominio público}          & \textbf{Proto} \\
\midrule
hermes-agent       & 127.0.0.1:8080 & \domain{hermes.el80.space}    & HTTP/WS \\
litellm-router     & 127.0.0.1:4000 & \domain{litellm.el80.space}   & HTTP    \\
grafana            & 127.0.0.1:3000 & \domain{grafana.el80.space}   & HTTP    \\
hermes-website     & 127.0.0.1:3001 & \domain{docs.el80.space}      & HTTP    \\
whisper-stt        & 127.0.0.1:9000 & — (interno)                   & HTTP    \\
prometheus         & 127.0.0.1:9090 & — (interno)                   & HTTP    \\
filebrowser        & 127.0.0.1:8095 & \domain{files.el80.space}     & HTTP    \\
syncthing (opt.)   & 127.0.0.1:8384 & \domain{syncthing.el80.space} & HTTP    \\
\bottomrule
\end{tabularx}
\end{table}

\section{Hardening de Seguridad en Contenedores}

Cada servicio aplica el principio de mínimo privilegio:

\begin{lstlisting}[language=yaml,caption=Configuración de seguridad base (docker-compose.yml)]
security_opt:
  - no-new-privileges:true
cap_drop:
  - ALL
cap_add:
  - CHOWN     # Solo cuando es necesario
  - SETGID
  - SETUID
tmpfs:
  - /app/tmp:size=64M,noexec,nosuid
\end{lstlisting}

Las excepciones justificadas son:
\begin{itemize}
  \item \service{cAdvisor}: requiere \code{privileged: true} para leer
    cgroups del kernel del host.
  \item \service{filebrowser}: corre como \code{user: "0:0"} para
    gestionar archivos de root con permisos completos.
  \item \service{autoheal}: monta \code{docker.sock} en modo read-only
    para detectar el estado de los contenedores.
\end{itemize}
